% PDF-derived LaTeX transcription. Original manuscript source was unavailable.
% Published PDF SHA-256: aa9dbe505497589c704154058617bb684e6ee875e84a053b6e87760fffbaed89
% Extracted text SHA-256: 7c586faf03129c766a3a94555d8df0a15edab7126e2efd3b8c67cbe21725c654
\documentclass[11pt]{article}
\usepackage[margin=1in]{geometry}
\usepackage{amsmath,amssymb}
\setlength{\emergencystretch}{3em}
\title{LUCAN-10B-Physics: A Sparse Foundation Model for Cross-Scale Physical Prediction}
\author{}
\date{}
\begin{document}
\maketitle
\noindent\textbf{PDF-derived LaTeX transcription.} This editable text follows the PDF extraction. Refer to the published PDF for figures, table layouts, and mathematical notation.

% Extracted chunk 1


\noindent SHODH AI RESEARCH\quad LUCAN-10B-Physics\par





LUCAN-10B-Physics: A Sparse Foundation
Model for Cross-Scale Physical Prediction


Shodh AI Technical Report
Shodh AI Research


\noindent 11.39B\quad 731,998\quad 14 of 25\par


\noindent total parameters, 9.8\% active on\quad tokens/s across eight H100s\quad primary continuum rows with\par


\noindent one token path\quad lowest error\par




\section*{Abstract}

LUCAN-10B-Physics is an 11.39B-parameter hybrid foundation model
that combines Mamba2 with 12 causal attention–mixture-of-experts layers,
each selecting one of 35 specialists alongside a shared expert. Only 9.8\%
of stored parameters lie on a single token path. Modality-specific stems
connect periodic crystal geometry, molecular fingerprints for biopharma-
ceutical properties, and continuum fields to one shared backbone.
Multi-donor sparse upcycling and FP32 router calibration lower latent val-
idation error by 10.5\%; the archived checkpoint exercises 412/420 routed
slots. LUCAN records the lowest error on 14 of 25 primary multi-step
comparisons against configured FNO and U-Net baselines on PDEBench,
with up to 1,000 test trajectories per task, and The Well. Simulation pre-
training reduces four-shot out-of-distribution error by 17.7\% for Cylinder
wake and 19.0\% for fluid–structure interaction.
The optional TCPI projection attains a maximum relative Sod mass resid-
ual of 6.57 \ensuremath{\times} 10\ensuremath{-}9 and 0.0\% Stefan positivity violations in paired replay.
Frozen-checkpoint scalar readouts improve over both matched descriptor
baselines on the full-split TDC AqSolDB corpus of 9,980 molecules and five-
fold Matbench Phonons. Eight H100s sustain 731,998 tokens/s of back-
bone throughput. These results establish measured operating points across
physical scales; constraint satisfaction, prediction accuracy, and hardware
throughput are evaluated separately.
An earlier 1,000-case Stefan workload gives a reported 71.1\ensuremath{\times} CPU wall-
time ratio at different achieved errors. On the first fixed Stefan test case,
the warmed v2 surrogate takes 23.64 ms, running 2.90\ensuremath{\times} faster than a
qualifying 512-cell H100 CUDA numerical solver (68.47 ms) and 10.45\ensuremath{\times}
faster than the CPU solver (246.99 ms).





\noindent Shodh AI Technical Report\quad 1\par


% Extracted chunk 2


\noindent SHODH AI RESEARCH\quad LUCAN-10B-Physics\par




\noindent SECTION 1\quad INTRODUCTION\par


A shared model across physical scales
Physical simulation supports engineering decisions across transport, fluid mechanics, phase
change, materials, and biopharmaceutical development. Numerical solvers encode governing
equations and boundary conditions explicitly, while learned surrogates trade a training in-
vestment for repeated prediction. Fourier Neural Operators (FNO) [1] and U-Net models [2]
provide strong task-specific baselines, but sharing representations across different equations,
geometries, and molecular inputs requires a broader interface.
Multiple Physics Pretraining (MPP) [3] studies autoregressive pretraining across hetero-
geneous physical systems. Poseidon [4] uses a multiscale operator transformer for PDE
transfer. WALRUS, introduced in November 2025 as arXiv:2511.15684, is a cross-domain
foundation model for continuum dynamics [5]. These systems motivate shared physical
representations. Our empirical rollout scorecard compares LUCAN with the fitted FNO
and U-Net baselines; it does not establish a measured ranking against Poseidon, MPP, or
WALRUS.
LUCAN-10B-Physics combines a causal Mamba2 recurrence [6] with 12 causal attention–
MoE blocks. Hybrid state-space, attention, and sparse-expert designs also appear in
Jamba [7]. We adapt that systems principle to a common width of 3,072 with task-specific
physical encoders and readouts. Each block contains 35 routed specialists and one shared
expert. The model stores 11.39 billion parameters, with 1,116,550,468 parameters on the
width-256 single-token path. The resulting 9.8\% structural fraction describes selective
execution, rather than batch memory or total attention cost.
Sparse initialization and routing. Sparse Upcycling reuses pretrained dense checkpoints
to initialize experts [8]; Drop-Upcycling introduces partial reinitialization to promote spe-
cialization [9]. Our multi-donor initialization, specialist adaptation, and train-fitted FP32
router calibration are evaluated both as an archived recipe and as a controlled three-seed
factorial. The recipe reduces pooled latent validation error from 1.20333 to 1.07667, while
the archived test audit reaches 412 of 420 routed layer–expert slots.

\noindent Continuum and experimental evidence.\quad We report matched horizons on\par

PDEBench [10] and The Well [11], and real-only versus simulation-then-real adaptation
on RealPDEBench, published at ICLR 2026 [12]. The 25 primary rollout rows use relative
L2 for PDEBench and VRMSE for The Well. LUCAN has the lowest unrounded error on
14 rows. Four-shot simulation pretraining lowers LUCAN’s Cylinder and FSI OOD errors
by 17.7\% and 19.0\%, respectively. Frozen-backbone readout controls measure the effect of
decoder choice.
Physical constraints and cross-modality reuse. TCPI imposes a prescribed global
inventory or admissible output bounds when the required physical information is available.
Paired replay measures its effect on identical decoded fields. The same frozen checkpoint
also supplies representations for Matbench crystal properties and TDC biopharmaceutical
and molecular properties. Matched raw-descriptor MLP and FP64 Ridge controls identify
where the residual representation helps and where simpler fits remain stronger.
Deployment scope. H100 measurements report backbone latency and data-parallel weak
scaling separately from complete surrogate execution. The earlier 1,000-case Stefan compar-
ison reports a 71.1\ensuremath{\times} CPU wall-time ratio at different achieved errors (Table 9). On the first
fixed Stefan test case, the warmed v2 surrogate takes 23.64 ms, running 2.90\ensuremath{\times} faster than
a qualifying 512-cell H100 CUDA numerical solver (68.47 ms) and 10.45\ensuremath{\times} faster than the
CPU solver (246.99 ms). Section 9 states the remaining timing and unified-pretraining
limitations; Appendix A records the dataset rosters and reproduction requirements.





\noindent Shodh AI Technical Report\quad 2\par


% Extracted chunk 3


\noindent SHODH AI RESEARCH\quad LUCAN-10B-Physics\par




\noindent SECTION 2\quad ARCHITECTURE\par


Cross-scale tokens and physical outputs


\section*{2.1 One token interface across physical scales}

Let um be a physical input, sm its modality stem, and Pm,t an optional patch or site indexing
operator. Learned lifts map native tokens to a common width:

\noindent xm,t = Lm (sm (Pm,t um )) ,\quad Lm : Rm \ensuremath{-}\ensuremath{\to} R3072 ,\quad m \ensuremath{\in} \{48, 192, 256\}.\par

Width 192 is used for crystal sites; width 256 is used for biopharmaceutical molecular
descriptors and continuum patches. Width 48 remains an additional supported interface.
A task mask identifies valid sites or patches and excludes padding from scalar pooling and
physical integrals.

\noindent Crystals\quad Biopharmaceutical /\quad Continuum fields\par

molecular properties

\noindent Crystal sites in periodic cells\quad Molecular graphs from SMILES Sampled fields and geometry\par


\noindent Species, distances and lattice\quad Radius-2 Morgan fingerprints  Local patches and physical\par


\noindent information\quad coordinates\par


\noindent Width-192 crystal tokens\quad Width-256 descriptor token\quad Width-256 grid tokens\par

Modality stems \ensuremath{-}\ensuremath{\to} learned lifts \ensuremath{-}\ensuremath{\to} shared 3,072-D LUCAN backbone

\noindent Masked pool + crystal head\quad Masked pool + property head\quad Multiscale field decoder +\par

optional TCPI

Figure 1: Scale ladder. Crystal, biopharmaceutical molecular-property, and continuum encoders
and readouts share the same 3,072-dimensional backbone.


\section*{2.2 Crystal distance–lattice stem}

For species Zi and Cartesian positions ri , the geometric encoder summarizes local composi-
tion and radial distances. For a periodic cell with lattice matrix A, define a minimum-image
displacement by
\ensuremath{\star}

\noindent nij \ensuremath{\in} arg min3 krj \ensuremath{-} ri + Ank2 ,\quad dij = krj \ensuremath{-} ri + Anij\par

\ensuremath{\star}
k2 .
n\ensuremath{\in}Z

For crystal sites, writing e(Z) for species features, \ensuremath{\rho}(d) for radial-basis features, and \ensuremath{\ell}(A)
for periodic-cell descriptors, the site pathway can be expressed as

\noindent [U+F8EE]\quad [U+F8F9]\par

X

\noindent ai = [U+F8F0]e(Zi ),\quad e(Zj ) \ensuremath{\otimes} \ensuremath{\rho}(dij ), \ensuremath{\ell}(A)[U+F8FB] ,  ssite\par


\noindent i\quad = E192 (ai ) \ensuremath{\in} R192 .\par

j6=i

The stem summarizes species, radial distances, and periodic-cell geometry for crystal-
property prediction. Minimum-image distances incorporate the crystal lattice without re-
quiring an additional force-field model.

\section*{2.3 Molecular graph and continuum stems}

A molecular graph G parsed from SMILES supplies a radius-2, 256-bit Morgan fingerprint, a
circular graph-derived descriptor [13]. With an optional train-fitted descriptor normalization
N , the molecular token is

\noindent f (G) = Morgan2,256 (G) \ensuremath{\in} \{0, 1\}256 ,\quad smol = E256 (N (f (G))) \ensuremath{\in} R256 .\par

This descriptor forms a single token for the molecular scalar pathway. It is a graph-derived
fingerprint rather than a claim that a new graph neural network was trained in these exper-
iments.
The continuum stem standardizes channels with training-split statistics, extracts local field
patches, and includes physical coordinates and patch metadata. In schematic form,

sgrid

\noindent t\quad = E256 ([N (Pt u), xt , \ensuremath{\mu}t ]) .\par




\noindent Shodh AI Technical Report\quad 3\par


% Extracted chunk 4


\noindent SHODH AI RESEARCH\quad LUCAN-10B-Physics\par



The task determines patch stride, channel layout, and valid-cell mask. The readout maps
contextual tokens back to the native prediction grid. In the reported decoders, raw fields
and coordinates also reach the readout alongside contextual backbone features.

\section*{2.4 Mamba2 recurrence and FP32 sparse routing}

The causal Mamba2 layer [6] combines a causal convolution with an input-dependent state-
space recurrence:
ht = Āt ht\ensuremath{-}1 + B̄t xt ,  zt = Ct ht + Dxt .
Twelve blocks follow, each with 24 causal attention heads and a sparse feed-forward module.
Given a block input v, router logits and probabilities are

\noindent g(v) = Wg v + bg\quad in FP32,\quad p(v) = softmax(g(v)),\par


\noindent r = arg max gj (v),\quad f (v) = pr (v)Er (v) + Eshared (v).\par

j\ensuremath{\in}\{1,...,35\}

The shared expert contributes on every token path. A batch can exercise several routed
specialists even though each token selects one per block. The recurrence has linear sequence-
length cost at fixed state width; the retained causal attention still has a quadratic token-pair
count. Measured wall time includes attention, dispatch, memory traffic, and kernel launches.

\section*{2.5 Multiscale spectral-dilated and scalar readouts}

The v1 spatial decoder combines Fourier modes, locally dilated convolutions, and low-rank
cross-attention to frozen backbone patches. Coordinates and an input residual retain native-
grid information. The staged v2 decoder adds a multiscale correction to a frozen v1 predic-

\noindent tion:\quad [U+0010]\quad [U+0011]\par

ybv2 = ybv1,frozen + \ensuremath{\Delta}multiscale xraw , \{h(3) , h(6) , h(9) , h(12) \} .
Signed 192-component projections mix the four backbone taps at each scale. Cross-attention
and FiLM condition a width-96, three-scale convolutional pyramid, with a full-resolution
Fourier branch retaining 16 or 24 modes. Readout selection uses the matched development-
validation metric, not the test score.
For a validity mask bt and last-layer token state ht , the scalar pathway uses masked pooling
and a residual readout:

\noindent P\quad [U+0010]\quad [U+0011]\par


\noindent bt ht\quad e .\par

e + \ensuremath{\Delta}\ensuremath{\theta} [h̄, d]

\noindent h̄ = Pt\quad ,\quad yb = fraw (d)\par

t bt

Here de denotes normalized raw descriptors. A feature-only head uses h̄ alone. Matched
descriptor-only MLP and Ridge controls retain the same raw information, target, and split;
the residual head adds contextual features and readout capacity.

\section*{2.6 TCPI: admissibility and exact inventory projection}

For a positive channel, softplus \ensuremath{\varphi}(ai ) = log(1+eai ) is a smooth option; the numerical control
uses \ensuremath{\varphi}(ai ) = max(0, ai ). Let qi = \ensuremath{\varphi}(ai ),PM \ensuremath{\geq} 0 be the input-derived inventory, vi > 0 the
cell volume, and V the valid-cell set. If j\ensuremath{\in}V vj qj > 0,

\noindent M qi\quad X\par


\noindent ybi = P\quad ,\quad vi ybi = M.\par


\noindent j\ensuremath{\in}V vj qj\quad i\ensuremath{\in}V\par


When M = 0, the positive-channel output is the zero field. A positive target with zero
normalizer requires an explicit runtime guard. The discrete ledger is checked after output
casting. For a signed periodic field, the affine correction
P
M \ensuremath{-} j\ensuremath{\in}V vj aj

\noindent ybi = ai +\quad P\par

j\ensuremath{\in}V vj

preserves the signed global integral. For Sod, the density inventory is initialized from the
input and advanced using known boundary fluxes. For Stefan, temperature is lower-clamped
and solid fraction uses
bi = min(1, max(0, ai )).
\ensuremath{\phi}



\noindent Shodh AI Technical Report\quad 4\par


% Extracted chunk 5


\noindent SHODH AI RESEARCH\quad LUCAN-10B-Physics\par



This phase map is piecewise differentiable, with unit derivative inside (0, 1) and zero outside.
Exact global inventory in arithmetic, numerical ledger residual, local PDE conservation,
and field accuracy are distinct properties. The paired experiments in Section 5 measure the
specified constraints and analytical-reference errors separately.


\noindent SECTION 3\quad MODEL SCALING\par


Upcycling and specialist routing
Multi-donor construction. Sparse Upcycling [8] initializes experts from pretrained check-
points, while Drop-Upcycling [9] motivates controlled reinitialization to encourage special-
ization. LUCAN’s five routed groups alternate between two donor checkpoints. A hidden-
channel permutation Pg transforms an expert’s two feed-forward matrices as

(W1 , W2 ) 7\ensuremath{-}\ensuremath{\to} (Pg W1 , W2 Pg> ),

preserving its function before Gaussian perturbation. Specialist adaptation and train-fitted
FP32 router calibration then refine the initialized model. The archived recipe lowers pooled
latent validation error from 1.20333 to 1.07667, a 10.5\% reduction.

\noindent A Historical native TEST routing\quad B Historical TEST\quad C Sealed checkpoint\par



\noindent Domain 1\quad 1.4\par


\noindent 15.0\quad 800\par

1.2




Within-layer pairs
12.5




Traffic share (\%)
Domain 2

\noindent 1.0\quad 600\par



Load CV
10.0

\noindent Domain 3\quad 0.8\par


\noindent 7.5\quad 400\par

0.6

\noindent Domain 4\quad 5.0\quad 0.4\par

200

\noindent 2.5\quad 0.2\par

Domain 5

\noindent 0.0\quad 0.0\quad 0\par


\noindent 1\quad 7\quad 14\quad 21\quad 28\quad 35\quad 1\quad 4\quad 8\quad 12\quad 0\quad 2\quad 4\quad 6\quad 8\par


\noindent Expert ID\quad Transformer layer\quad FC1 + FC2 cosine\quad 1e\ensuremath{-}3\par





Figure 2: Routing, traffic and expert-weight measurements on the archived five-domain evaluation.
The checkpoint uses 412/420 routed layer–expert slots. The maximum absolute unaligned expert-
weight cosine is 0.008469 across 7,140 pairs. CV denotes traffic coefficient of variation.
Table 1: Archived multi-donor upcycling, specialist adaptation and FP32 router calibration com-
pared with an uncalibrated clone. Error is pooled latent relative L2 . This recipe comparison is
distinct from the matched factorial below.


\noindent Measurement\quad Uncalibrated clone Calibrated LUCAN-10B\par


\noindent Validation error\quad 1.20333\quad 1.07667\par


\noindent Active validation slots\quad 81/420\quad 310/420\par


\noindent Mean validation layer CV\quad 3.21\quad 1.77\par


\noindent Active slots on archived test panel\quad —\quad 412/420\par


\noindent Maximum absolute unaligned weight cosine\quad 1.000000\quad 0.008469\par



\section*{3.1 Three-seed donor-by-calibration factorial}

The controlled factorial holds specialist adaptation, token exposure, and Gaussian pertur-
bation scale \ensuremath{\sigma} = 0.005 fixed. The single-donor arm repeats the same donor checkpoint; the
multi-donor arm uses both. The calibration contrast toggles a router bias fitted on TRAIN.
Validation contains 12 cases; the sealed five-domain test split contains 10 cases. Table 2
reports arm means and the available paired 95\% confidence intervals for the effects.





\noindent Shodh AI Technical Report\quad 5\par


% Extracted chunk 6


\noindent SHODH AI RESEARCH\quad LUCAN-10B-Physics\par



Table 2: Three-seed factorial with matched specialist adaptation. Upper panel: mean errors and
mean active test slots. Lower panel: paired TEST contrasts with their reported 95\% CIs. A
negative contrast favors the added component; intervals describe contrasts, rather than per-arm
error uncertainty.

Donors FP32 calibration VAL error TEST error TEST slots / 420

\noindent Single\quad Off\quad 1.20839\quad 1.17993\quad 193.3\par


\noindent Single\quad On\quad 1.20708\quad 1.17747\quad 394.7\par


\noindent Multi\quad Off\quad 1.20430\quad 1.17552\quad 192.3\par


\noindent Multi\quad On\quad 1.20488\quad 1.17518\quad 395.0\par



\noindent Paired TEST contrast\quad Mean difference\quad 95\% CI\par


\noindent Multi minus single donor, averaged over calibration\quad \ensuremath{-}0.003350 [\ensuremath{-}0.005405, \ensuremath{-}0.002826]\par


\noindent Calibration on minus off, averaged over donors\quad \ensuremath{-}0.001401 [\ensuremath{-}0.002264, \ensuremath{-}0.000724]\par


\noindent Donor–calibration interaction\quad +0.002111 [+0.001210, +0.005485]\par

The positive interaction indicates a smaller calibration gain with multi-donor initialization.
The validation calibration interval includes zero. Within the multi-donor test arm, cali-
bration increases observed slot coverage from 45.8\% to 94.0\%. The archived checkpoint’s
412-slot audit and the factorial’s 395.0 mean slots are different evaluations and should not
be substituted for each other.

\section*{3.2 Error aggregation and expert diversity}

The pooled latent score weights domains by target-squared magnitude. One domain supplies
89.6\% of the denominator on the ten-case test panel. Equal-domain mean contrasts remain
negative: \ensuremath{-}0.003462 for donors and \ensuremath{-}0.001777 for calibration, with negative point contrasts
in each domain. The 10.5\% archived recipe improvement should therefore be read alongside
the controlled factorial, rather than attributed to calibration alone.
Unaligned weight cosine is sensitive to the construction permutations and does not estab-
lish functional diversity by itself. After reversing those permutations, the mean seed-level
median cosine changes from 0.002452 to 0.002345 between single- and multi-donor experts.
Functional-output cosine changes from 0.6224 to 0.5896 and directional-Jacobian cosine
from 0.2138 to 0.2014. Within each arm, experts share pre-router probe inputs; cross-arm
activations come from each model’s own held-out pass.


\noindent SECTION 4\quad CONTINUUM PREDICTION\par


Multi-step rollout evaluation


\section*{4.1 Matched horizons and trajectory partitions}

PDEBench [10] uses 1,000 test trajectories for 1D Advection, 1D Burgers, 1D Reaction–
Diffusion, and 2D Darcy Flow; 2D Reaction–Diffusion uses 100. Darcy is a static coefficient-
to-pressure map. Dynamic tasks begin from the same observed frame and recursively feed
predictions into later steps. The Well [11] retains its published whole-trajectory partitions
for Gray–Scott, Active Matter, Shear Flow, and Rayleigh–Bénard. Appendix A gives the
complete rosters.
For a prediction at horizon t, relative L2 compares the field error norm with the target
norm. The Well uses per-channel variance-normalized RMSE,
s
MSE(b  yn,c,t , yn,c,t )

\noindent VRMSEn,c,t =\quad ,\par

Vari (yn,i,c,t ) + 10\ensuremath{-}7

aggregated across channels and trajectories. Active Matter relative-L2 rows are supplemen-
tary. The 25 primary rows comprise the 13 PDEBench rows and 12 The Well VRMSE rows;
they do not include those three supplementary rows.
Validation-selected readouts. For The Well, the spatial objective combines variance-
normalized RMSE with relative L2 and a one- to three-step pushforward curriculum. A
staged v2 candidate replaces v1 only if it improves the matched validation metric: first-step



\noindent Shodh AI Technical Report\quad 6\par


% Extracted chunk 7


\noindent SHODH AI RESEARCH\quad LUCAN-10B-Physics\par



mean VRMSE for The Well and mean relative L2 for PDEBench and RealPDEBench. This
gate selected 21 of 23 candidates across rollout and transfer experiments. The Burgers and
Darcy rows retain v1, with no v2 candidates in that sweep; the other LUCAN rollout rows
use v2.

\section*{4.2 Field comparisons and rollout behavior}

LUCAN records the lowest error in 14 of 25 primary comparisons. Relative to FNO,
static Darcy error is approximately 8.5\ensuremath{\times} lower; Advection at step 10 is 8.5\ensuremath{\times} lower; and
Rayleigh–Bénard step-10 VRMSE is 3.75\ensuremath{\times} lower. These are row-specific ratios, rather than
an aggregate benchmark speed or accuracy factor.
Gray–Scott and Rayleigh–Bénard favor LUCAN at longer reported horizons, while U-Net
remains strongest on the primary Active Matter rows. Shear Flow favors LUCAN at steps
1 and 5 and FNO at step 10. Reporting all horizons and baselines preserves these operating
differences rather than reducing the study to the win count.
Table 3: Continuum prediction on PDEBench and The Well at t = 1, 5, 10, with one static
prediction for Darcy. N counts test trajectories. Lower is better; bold marks the lowest unrounded
error among the three configured methods.


\noindent Dataset\quad N\quad t Measure LUCAN-10B\quad FNO\quad U-Net\par

PDEBench

\noindent 1D Advection\quad 1000\quad 1\quad Rel. L2\quad 0.00067 0.00329  0.00061\par


\noindent 5\quad Rel. L2\quad 0.00185 0.01168\quad 0.00451\par


\noindent 10\quad Rel. L2\quad 0.00333 0.02848\quad 0.02190\par


\noindent 1D Burgers\quad 1000\quad 1\quad Rel. L2\quad 0.00401 0.00508\quad 0.00460\par


\noindent 5\quad Rel. L2\quad 0.05303 0.05682\quad 10.88610\par


\noindent 10\quad Rel. L2\quad 0.09942 0.12530 8.31 \ensuremath{\times} 108\par


\noindent 1D Reaction–Diffusion 1000\quad 1\quad Rel. L2\quad 0.00014 0.00014\quad 0.07650\par


\noindent 5\quad Rel. L2\quad 0.00094 0.00068\quad 1.20228\par


\noindent 10\quad Rel. L2\quad 0.00138 0.00115  19.33048\par


\noindent 2D Darcy Flow\quad 1000\quad 1\quad Rel. L2\quad 0.02951 0.25136\quad 0.51650\par


\noindent 2D Reaction–Diffusion 100\quad 1\quad Rel. L2\quad 0.96580 0.97531  0.96537\par


\noindent 5\quad Rel. L2\quad 0.97628 1.03273  0.97377\par


\noindent 10\quad Rel. L2\quad 0.97525 1.05820\quad 0.98600\par

The Well

\noindent Gray–Scott\quad 120\quad 1\quad VRMSE\quad 0.18929\quad 0.22863  0.18314\par


\noindent 5\quad VRMSE\quad 0.58317\quad 1.03718\quad 0.91264\par


\noindent 10\quad VRMSE\quad 1.33598\quad 1.96542\quad 1.94491\par


\noindent Active Matter\quad 26\quad 1\quad VRMSE\quad 3.23628\quad 5.33087  2.54440\par


\noindent 5\quad VRMSE\quad 14.58457\quad 55.67803 11.90089\par


\noindent 10\quad VRMSE\quad 13.08445\quad 29.87972  7.51500\par


\noindent 26\quad 1\quad Rel. L2\quad 0.12295\quad 0.16503\quad 0.20068\par


\noindent 5\quad Rel. L2\quad 0.91517\quad 2.53734  0.83014\par


\noindent 10\quad Rel. L2\quad 1.22464\quad 2.16258  1.00621\par


\noindent Shear Flow\quad 112\quad 1\quad VRMSE\quad 0.01636\quad 0.03080\quad 0.23875\par


\noindent 5\quad VRMSE\quad 0.17242\quad 0.22500\quad 0.86179\par


\noindent 10\quad VRMSE\quad 0.52795\quad 0.45843\quad 6.87916\par


\noindent Rayleigh–Bénard\quad 175\quad 1\quad VRMSE\quad 0.43313\quad 1.87671\quad 3.02995\par


\noindent 5\quad VRMSE\quad 2.46888\quad 12.73057  51.85465\par


\noindent 10\quad VRMSE\quad 5.88615\quad 22.09398 178.92939\par

The 25 primary comparisons use relative L2 for PDEBench and VRMSE for The Well. The three
Active Matter relative-L2 rows report a different weighting of field error and are supplementary.





\noindent Shodh AI Technical Report\quad 7\par


% Extracted chunk 8


\noindent SHODH AI RESEARCH\quad LUCAN-10B-Physics\par




\noindent Reference\quad LUCAN-10B\quad FNO\quad U-Net\par





Darcy




Reaction–diffusion




Cylinder wake




1.0

0.8


Sod density
0.6

0.4

0.2



\noindent 0\quad 1\quad 0\quad 1\quad 0\quad 1\quad 0\quad 1\par


1.0

0.8




Stefan phase
0.6

0.4

0.2

0.0

\noindent 0\quad 1\quad 0\quad 1\quad 0\quad 1\quad 0\quad 1\par





Figure 3a: Matched reference and predicted fields for Darcy, reaction–diffusion, Cylinder wake,
Sod density, and Stefan phase fraction. Each row uses the same field or case across models.

\noindent Gray–Scott\quad Shear Flow\quad Cylinder · OOD\par


\section*{2.0                                                                                                                                                                                           LUCAN-10B-Physics}


\noindent 0.12\quad FNO\par





VRMSE (log scale)
U-Net
1.5



Relative L2

\noindent 100\quad 0.11\par



\noindent VRMSE\quad 1.0\par

0.10
10\ensuremath{-}1

\noindent 0.5\quad 0.09\par




\noindent 1\quad 3\quad 5\quad 7\quad 10\quad 1\quad 3\quad 5\quad 7\quad 10\quad 4\quad 8\quad 16\quad 32\par


\noindent Forecast step\quad Forecast step\quad Real training trajectories\par


\noindent Active Matter\quad Rayleigh–Bénard\quad FSI · OOD\par

60

\noindent 102\quad 0.085\par





VRMSE (log scale)
0.080


Relative L2
40

VRMSE

\noindent 101\quad 0.075\par



\noindent 20\quad 0.070\par



\noindent 100\quad 0.065\par


0

\noindent 1\quad 3\quad 5\quad 7\quad 10\quad 1\quad 3\quad 5\quad 7\quad 10\quad 4\quad 8\quad 16\quad 32\par


\noindent Forecast step\quad Forecast step\quad Real training trajectories\par





Figure 3b: Multi-step The Well rollouts and OOD transfer after simulation pretraining for Cylin-
der and FSI. Rollout panels show mean VRMSE; transfer panels show mean trajectory relative L2
against the number of real training trajectories. The original measured curves are retained.





\noindent Shodh AI Technical Report\quad 8\par


% Extracted chunk 9


\noindent SHODH AI RESEARCH\quad LUCAN-10B-Physics\par




\noindent SECTION 5\quad EMPIRICAL CONTROLS\par


Readouts and physical conservation


\section*{5.1 Frozen-backbone readout ablation}

The original readout comparison holds the frozen LUCAN backbone and physical inputs
fixed while fitting a one-layer linear probe, a two-layer MLP, or a two-layer spectral-dilated
decoder. Table 4 preserves both field errors and the applicable absolute mass-ledger resid-
uals. These original readout values precede the staged validation-selected v2 scorecard in
Table 3; the two tables measure different readout protocols.
Table 4: Readout ablation on the same frozen backbone. Relative L2 and absolute signed mass-
ledger differences are evaluated on matched physical grids. Dashes indicate that no applicable
ledger is reported. Bold marks the lowest unrounded mean within each metric block.


\noindent Relative L2\quad Absolute mass residual\par


\noindent Dataset\quad N\quad t\quad Linear\quad MLP Spectral Linear MLP\quad Spectral\par


\noindent 1D Advection\quad 1 0.06352 0.05718 0.00078 0.0152 0.0004 4.41 \ensuremath{\times} 10\ensuremath{-}5\par

1000

\noindent 5 0.30958 0.27710 0.00247 0.0772 0.0018\quad 0.0002\par


\noindent 10 0.58900 0.52053 0.00457 0.1579 0.0036\quad 0.0005\par


\noindent 1D Burgers\quad 1000 1 0.09792 0.08253 0.00401 0.0262 0.0028\quad 0.0003\par


\noindent 5 0.40724 0.33124 0.05303 0.1137 0.0128\quad 0.0032\par


\noindent 10 0.64361 0.51210 0.09942 0.1901 0.0240\quad 0.0084\par


\noindent 1D Reaction–Diffusion 1000 1 0.13120 0.07441 0.00043\quad —\quad —\quad —\par


\noindent 5 0.05496 0.06309 0.00177\quad —\quad —\quad —\par


\noindent 10 0.09831 0.05644 0.00380\quad —\quad —\quad —\par


\noindent 2D Darcy Flow\quad 1000 1 0.50575 0.49913 0.02951\quad —\quad —\quad —\par


\noindent 2D Reaction–Diffusion 100 1 1.04557 0.98065 0.96609\quad —\quad —\quad —\par


\noindent 5 1.08538 1.04411 0.99118\quad —\quad —\quad —\par


\noindent 10 1.32491 1.14829 1.03735\quad —\quad —\quad —\par


\noindent Euler Sod\quad 64 1 0.36738 0.26095 0.05930\quad —\quad —\quad —\par


\noindent Stefan phase change\quad 64 1 0.01163 0.01003 0.00019\quad —\quad —\quad —\par

At the first prediction, the spectral decoder reduces relative L2 by approximately 81\ensuremath{\times} on
Advection, 24\ensuremath{\times} on Burgers, 303\ensuremath{\times} on 1D Reaction–Diffusion, and 17\ensuremath{\times} on Darcy relative
to the linear probe. The corresponding Sod and Stefan ratios are approximately 6.2\ensuremath{\times} and
60\ensuremath{\times}. The smaller change on 2D Reaction–Diffusion is retained in the full table. These
controls establish the importance of the readout at fixed backbone features; they do not by
themselves isolate the value of backbone pretraining.

\section*{5.2 Discontinuities, positivity and exact inventories}

Table 5 compares neural operators against analytical references on 64 held-out cases per
task. Table 6 uses a separate 16-case cohort per task and switches TCPI off and on for
the same decoded fields. This pairing isolates the projection from changes in training data,
backbone, or decoder.
Table 5: Held-out analytical neural-operator comparisons, N = 64 per task. Lower field error and
positivity-violation rate are better.


\noindent Task\quad Metric\quad LUCAN-10B\quad FNO\quad U-Net\par


\noindent Euler Sod\quad Mean Rel. L2\quad 0.059302\quad 0.065957 0.071271\par


\noindent Positivity violations (\%)\quad 0.00\quad 0.00\quad 0.00\par


\noindent Stefan phase change Mean Rel. L2\quad 0.000192 7.92 \ensuremath{\times} 10\ensuremath{-}5 0.000274\par


\noindent Positivity violations (\%)\quad 2.57\quad 20.90\quad 19.74\par






\noindent Shodh AI Technical Report\quad 9\par


% Extracted chunk 10


\noindent SHODH AI RESEARCH\quad LUCAN-10B-Physics\par



Table 6: Matched TCPI replay, N = 16 per task. Raw and constrained predictions share the same
decoded cases. A dash indicates that the displayed mass-residual ledger does not apply.


\noindent Task\quad Metric\quad LUCAN raw LUCAN + TCPI\par


\noindent Euler Sod\quad Pooled Rel. L2\quad 0.11526\quad 0.08923\par


\noindent Positivity violations (\%)\quad 0.00\quad 0.00\par


\noindent Max. rel. mass residual\quad 0.15622\quad 6.57 \ensuremath{\times} 10\ensuremath{-}9\par


\noindent Stefan phase change Pooled Rel. L2\quad 0.00957\quad 0.00451\par


\noindent Positivity violations (\%)\quad 37.16\quad 0.00\par


\noindent Max. rel. mass residual\quad —\quad —\par

On the 64-case Sod comparison, the spectral readout reaches 0.059302 mean relative L2 ,
versus 0.065957 for FNO and 0.071271 for U-Net. In the paired replay, TCPI lowers pooled
Sod error from 0.11526 to 0.08923 and the maximum relative mass residual from 0.15622 to
6.57 \ensuremath{\times} 10\ensuremath{-}9 , with zero positivity violations. For Stefan, paired positivity violations fall from
37.16\% to 0.00\% and pooled error from 0.00957 to 0.00451. The separate 64-case Stefan
table retains FNO’s lower field error and the raw spectral decoder’s 2.57\% violation rate.

\noindent Sod shock tube\quad Stefan phase change\par


\section*{1.0                                                                                                                 Analytical reference}

6
LUCAN raw
LUCAN constrained
0.8
4




Solid fraction
Density
0.6
2

0.4
0

0.2
\ensuremath{-}2


\noindent 0.00\quad 0.25\quad 0.50\quad 0.75\quad 1.00\quad 0\quad 5\quad 10\quad 15\quad 20\par


\noindent Position (normalized)\quad Position (mm)\par




Figure 4: Sod density and Stefan phase fraction on representative paired-replay cases. Analytical
references, raw LUCAN outputs, and constrained outputs share the same coordinates.
The separate 1,000-case Stefan workload appears in Table 9 in Section 8. Its CPU and H100
operating points have different achieved errors and are distinct from both analytical cohorts
above.


\noindent SECTION 6\quad EXPERIMENTAL TRANSFER\par


Simulation helps when real data is scarce
RealPDEBench provides paired simulated and experimental velocity fields for Cylinder wake
and fluid–structure interaction (FSI) [12]. We compare fitting on real data alone with
simulation pretraining followed by experimental adaptation. Both protocols use the frozen
11.39B backbone and the same experimental trajectory budget.

\section*{6.1 Trajectory budgets, windows and OOD scoring}

A shot is a distinct experimental training trajectory, with k \ensuremath{\in} \{4, 8, 16, 32\}. The full test
rosters cover 4,827 windows from 92 Cylinder trajectories and 2,731 windows from
45 FSI trajectories. OOD scores average the held-out trajectories: 10 for Cylinder and
12 for FSI. Window counts describe the full test panels, rather than the number of OOD
trajectories or independent training shots.
Development excludes test overlap, deduplicates FSI validation coordinates, and preserves
published test-row multiplicity. A fixed stem supplies cached v1 features; staged v2 fits a
residual on frozen v1. Validation gating retains v1 for real-only Cylinder at four and eight
shots. Both protocols use the same mean-trajectory relative-L2 metric and shot budget;
Appendix A gives the window-sampling rules.


\noindent 17.7\%\quad 19.0\%\par


\noindent Cylinder wake error reduction\quad FSI error reduction\par


\noindent 0.13621 to 0.11207\quad 0.08564 to 0.06935\par




\noindent Shodh AI Technical Report\quad 10\par


% Extracted chunk 11


\noindent SHODH AI RESEARCH\quad LUCAN-10B-Physics\par



Table 7: RealPDEBench OOD mean-trajectory relative L2 . One shot equals one experimental
training trajectory. Bold marks the lowest unrounded error among the three models within each
protocol and shot budget.


\noindent Real-only\quad Simulation then real\par


\noindent Scenario\quad k LUCAN-10B\quad FNO\quad U-Net LUCAN-10B\quad FNO\quad U-Net\par


\noindent Cylinder\quad 4\quad 0.13621\quad 0.14667 0.11105\quad 0.11207\quad 0.12469 0.10852\par


\noindent 8\quad 0.10929\quad 0.11746 0.10087\quad 0.10132\quad 0.10409 0.10062\par


\noindent 16\quad 0.08623\quad 0.09352 0.08786\quad 0.08581\quad 0.09495 0.09047\par


\noindent 32\quad 0.08236\quad 0.09030 0.08652\quad 0.08294\quad 0.09091 0.08818\par


\noindent FSI\quad 4\quad 0.08564\quad 0.10382\quad 0.07100\quad 0.06935\quad 0.08748\quad 0.06648\par


\noindent 8\quad 0.07374\quad 0.09202\quad 0.06647\quad 0.06606\quad 0.08338\quad 0.06431\par


\noindent 16\quad 0.06819\quad 0.08371\quad 0.06382\quad 0.06463\quad 0.08123\quad 0.06323\par


\noindent 32\quad 0.06438\quad 0.07925\quad 0.06236\quad 0.06263\quad 0.07763\quad 0.06159\par



\section*{6.2 Effect of simulation pretraining}

At four shots, LUCAN’s Cylinder error falls from 0.13621 to 0.11207 (17.7\%), and FSI
error from 0.08564 to 0.06935 (19.0\%). These compare the same model’s real-only and
simulation-then-real fits.
LUCAN leads Cylinder at 16 and 32 shots under both protocols, four of the 16 scenario–
protocol–shot comparisons. U-Net leads every displayed FSI row. Simulation-pretraining
gains are distinct from baseline win rates.


\noindent SECTION 7\quad CROSS-MODALITY TRANSFER\par


Crystals and biopharmaceutical properties
The same frozen checkpoint used for continuum prediction supplies representations for
Matbench crystal properties [14] and Therapeutics Data Commons biopharmaceutical and
molecular properties [15]. The periodic crystal-site stem has width 192, and the molecular
graph-derived descriptor stem has width 256. Task-specific encoders and lightweight scalar
heads adapt to each input and target; the comparison holds descriptor inputs and splits
fixed across the LUCAN residual head, raw MLP, and FP64 Ridge controls.

\section*{7.1 Scalar protocol and aggregation}

For a training set Dtr , descriptor normalization is fitted on training inputs and applied
unchanged to validation and test. Let den = Ntr (dn ), let h̄n be the masked pooled frozen-
backbone state, and let yn be the target in its native units. The three prediction families
are
ybnMLP = f\ensuremath{\psi} (den ),
ybRidge = w> den + b,

\noindent n\quad \ensuremath{\lambda}\par


ybnLUCAN = fraw (den ) + \ensuremath{\Delta}\ensuremath{\theta} ([h̄n , den ]).
\ensuremath{-}1 P

\noindent The per-split score is MAE = Ntest\quad n |b\par

yn \ensuremath{-} yn |. Model selection uses validation MAE.
The development panel uses one Matbench fold and one TDC scaffold split. The full-split
panel uses all five Matbench folds and three TDC scaffold seeds. Matbench and TDC
values average fold or seed MAEs. TDC test rosters can overlap across seeds, so they are
not three disjoint test populations. The reported AqSolDB corpus contains 9,980 molecules
before scaffold partitioning; that corpus size is not the held-out sample count for every seed.





\noindent Shodh AI Technical Report\quad 11\par


% Extracted chunk 12


\noindent SHODH AI RESEARCH\quad LUCAN-10B-Physics\par



Table 8: Scalar test MAE for LUCAN-10B and matched descriptor-only baselines. Development-
panel and full-split results have different coverage. Bold marks the lowest point estimate among
the three methods within each panel; native target units are retained.


\noindent Development-panel MAE\quad Full-split MAE\par


\noindent Dataset\quad Unit\quad LUCAN MLP\quad Ridge LUCAN\quad MLP Ridge\par


\noindent Matbench JDFT2D\quad meV/atom\quad 46.55\quad 50.99\quad 92.05\quad 47.57 47.07\quad 64.28\par


\noindent Matbench Phonons\quad cm\ensuremath{-}1\quad 70.50\quad 74.58\quad 86.39\quad 77.24 78.35\quad 95.22\par


\noindent TDC AqSolDB\quad log mol/L\quad 1.494\quad 1.509\quad 1.550\quad 1.429 1.442\quad 1.631\par


\noindent TDC Lipophilicity\quad log D\quad 0.748\quad 0.751\quad 0.821\quad 0.778 0.763\quad 0.814\par

Across the retained crystal and biopharmaceutical tasks, LUCAN-10B improves over both
descriptor-only baselines on AqSolDB (1.429 versus 1.442 for MLP and 1.631 for Ridge) and
Matbench Phonons (77.24 versus 78.35 and 95.22). The other full-split results are closely
comparable with the best descriptor baseline: JDFT2D is 47.57 versus 47.07 for MLP, and
Lipophilicity is 0.778 versus 0.763. These are point-estimate comparisons; the table does
not establish statistical significance for the small MAE differences.

\section*{7.2 Heterogeneous scaffolds and multi-species crystals}

AqSolDB spans heterogeneous molecular scaffolds, and Phonons spans different crystal com-
positions and structures. The frozen nonlinear representation and residual readout can en-
code interactions that are less accessible to one global linear descriptor map. This interpre-
tation is consistent with the lower AqSolDB and Phonons MAEs and the close comparisons
on JDFT2D and Lipophilicity.
This interpretation is a hypothesis about representation and bias–variance tradeoffs, rather
than a causal identification from these tables. The residual pathway also adds head capacity.
The smaller development study was exploratory: its residual design had been informed
by earlier test results before validation-based checkpoint selection. The broader full-split
evaluation increases coverage, but independent blinded confirmation remains separate from
the reported point estimates.


\noindent SECTION 8\quad GPU DEPLOYMENT\par


Measured H100 latency and throughput


\section*{8.1 Single-device latency and measurement boundaries}

The synthetic-token profile uses NVIDIA H100 80 GB HBM3, BF16 width-256 inputs, FP32
routing, and PyTorch 2.10.0+cu128. Each single-device configuration uses ten synchronized
timed forwards after warm-up. Across the reported configurations, the stated latency end-
points are 18.1–47.95 ms; the 18.1 ms value belongs to a separate earlier setting and is
not assigned a sequence length in the controlled series below.
The directly comparable batch-one sequence sweep measures 26.43, 27.77, 30.54, 35.34,
and 47.95 ms at 256, 512, 1,024, 2,048, and 4,096 tokens. A 16\ensuremath{\times} sequence-length increase
therefore costs 1.81\ensuremath{\times} measured wall time in that sweep. At 256 tokens and batch 64, the
median is 100.18 ms. These are backbone measurements, rather than complete encoder–
decoder workload timings. The model retains causal attention, so the sequence sweep is not
evidence of an entirely linear-cost architecture.

\section*{8.2 Data-parallel weak scaling}

With 256 tokens per sequence and batch 32 per data-parallel replica, summed throughput
is 126,181, 238,381, 366,472, and 731,998 tokens/s on one, two, four, and eight H100s.
Eight GPUs achieve 5.80\ensuremath{\times} the one-GPU rate and 72.5\% parallel efficiency. Peak allocated
memory is 23.66 GB per device in the eight-GPU run.





\noindent Shodh AI Technical Report\quad 12\par


% Extracted chunk 13


\noindent SHODH AI RESEARCH\quad LUCAN-10B-Physics\par






\noindent A Single-H100 latency\quad B Data-parallel weak scaling\par



\noindent Batch one\quad 5.80\ensuremath{\times} gain; 72.5\% efficiency\par





Throughput (1,000 tokens/s)

\noindent 50\quad 1000\par


\noindent Earlier setting: 18.1 ms*\quad Linear ideal\par





Median latency (ms)

\noindent 45\quad 800\quad Measured\par


40
600
35
400
30
200
25


\noindent 20\quad 0\par


\noindent 6\quad 2\quad 4\quad 8\quad 6\quad 1\quad 2\quad 4\quad 8\par


\noindent 25\quad 51\quad 1\quad ,02\quad 2\quad ,04\quad 4\quad ,09\par

H100 GPUs
Tokens per sequence




Figure 5: H100 deployment measurements. (A) Controlled batch-one backbone latency against
sequence length; the separate 18.1 ms earlier endpoint is annotated without placing it on that curve.
(B) Data-parallel weak scaling at batch 32 per GPU, with linear ideal scaling shown for reference.
Only measured configurations are plotted.

\section*{8.3 Solver crossover and evaluation scope}

On the first fixed Stefan test case, the warmed v2 surrogate takes 23.64 ms, running 2.90\ensuremath{\times}
faster than a qualifying 512-cell H100 CUDA numerical solver (68.47 ms) and 10.45\ensuremath{\times}
faster than the CPU solver (246.99 ms). At batch 1,000, coarser numerical solvers qualify
and run faster.
The earlier 1,000-case workload in Table 9 takes 2.97 s for the H100 surrogate and 210.94
s for the 512-cell CPU enthalpy solver. The corresponding pooled relative-L2 errors are
0.00279 and 1.67 \ensuremath{\times} 10\ensuremath{-}5 . The 71.1\ensuremath{\times} wall-time ratio describes these two operating points.
It cannot be interpreted as an accuracy-matched solver speedup or extrapolated to different
geometries, hardware, batch sizes, or numerical resolutions.
Table 9: Earlier solver workloads at different accuracy and an accuracy-qualified batch-one Stefan
comparison. Panel (a) reports total wall time and pooled relative L2 for the earlier 1,000-case
workload. Panel (b) reports latency on the first fixed Stefan test case; numerical solvers use 512
cells. Time ratios in panel (b) divide numerical-solver latency by warmed v2-surrogate latency.

(a) Earlier 1,000-case workload: different achieved errors


\noindent Method\quad Hardware / grid\quad Wall time (s) Pooled relative L2\par


\noindent LUCAN surrogate One NVIDIA H100\quad 2.97\quad 0.00279\par


\noindent Enthalpy solver CPU, 512 cells\quad 210.94\quad 1.67 \ensuremath{\times} 10\ensuremath{-}5\par


(b) First fixed Stefan test case: batch one


\noindent Method\quad Hardware / grid\quad Latency (ms) Time / surrogate\par


\noindent Warmed v2 surrogate\quad One NVIDIA H100\quad 23.64\quad 1.00\ensuremath{\times}\par


\noindent Qualifying CUDA solver H100, 512 cells\quad 68.47\quad 2.90\ensuremath{\times}\par


\noindent CPU solver\quad CPU, 512 cells\quad 246.99\quad 10.45\ensuremath{\times}\par

Backbone token throughput, complete surrogate latency, and numerical-solver wall time
answer different deployment questions. Accuracy-matched comparisons require a common
reference, physical inputs, output grid, forecast interval, and timing boundary, with numer-
ical resolution chosen independently to meet the same error and constraint criteria.





\noindent Shodh AI Technical Report\quad 13\par


% Extracted chunk 14


\noindent SHODH AI RESEARCH\quad LUCAN-10B-Physics\par




\noindent SECTION 9\quad LIMITATIONS\par


Scope and remaining scientific boundaries
Non-equivalent solver timings. The earlier 1,000-case Stefan CPU and H100 measure-
ments use different hardware and achieve different errors. Their wall-time ratio is a workload
observation, rather than an accuracy-matched acceleration claim. Likewise, single-device se-
quence sweeps and data-parallel weak scaling have different execution boundaries. TCPI’s
global inventory or admissibility guarantee does not establish local PDE consistency, point-
wise accuracy, or long-horizon stability.
Future unified cross-modality pretraining. The shared checkpoint is reused through
separate modality stems and task heads; the present experiments do not demonstrate one
joint crystal–molecular–continuum pretraining objective or a universally aligned tokenizer.
Future work should jointly train and evaluate those interfaces on disjoint molecular, crystal,
and continuum systems, with independent test custody. Earlier panel results informed
some readout refinements, so independent blinded evaluation remains necessary to confirm
transfer beyond the reported benchmark and scalar coverage.





\noindent Shodh AI Technical Report\quad 14\par


% Extracted chunk 15


\noindent SHODH AI RESEARCH\quad LUCAN-10B-Physics\par



APPENDIX A

Dataset rosters and reproducibility

A.1 Published trajectory rosters
Table A1 records the trajectory populations used by the continuum experiments. Counts are
trajectories, rather than patches, training pairs, or sampled windows. The DaRUS entries
are public source-file identifiers for the PDEBench data, not internal execution identifiers.
Table A1: Continuum trajectory rosters. Train and VAL form development populations; TEST
is the held-out roster. Dynamic prediction and static Darcy evaluation use the task definitions in
Section 4.


\noindent Task\quad Source pin\quad Train\quad VAL\quad TEST\par


\noindent 1D Advection\quad DaRUS 255674\quad 8100\quad 900\quad 1000\par


\noindent 1D Burgers\quad DaRUS 281363\quad 8100\quad 900\quad 1000\par


\noindent 2D Darcy\quad DaRUS 133219\quad 8100\quad 900\quad 1000\par


\noindent 1D Reaction–Diffusion\quad DaRUS 133181\quad 8100\quad 900\quad 1000\par


\noindent 2D Reaction–Diffusion\quad DaRUS 133017\quad 810\quad 90\quad 100\par


\noindent Active Matter\quad The Well a\quad 175\quad 24\quad 26\par


\noindent Gray–Scott\quad The Well a\quad 960\quad 120\quad 120\par


\noindent Rayleigh–Bénard\quad The Well a\quad 1400\quad 175\quad 175\par


\noindent Shear Flow\quad The Well a\quad 896\quad 112\quad 112\par

a
The Well uses its published whole-trajectory partitions [11]. Numerical reproduction must pin
the corresponding source files and their byte-level identities.
PDEBench keeps the published outer test roster and ranks development trajectories by
SHA-256 for validation. The Well follows upstream train, validation, and test partitions.
Dynamic-grid fits use four equally spaced training starts per trajectory and validate from
start zero. Darcy uses one static input–target map. Normalization statistics are fitted on
training inputs only.
A.2 Experimental and scalar partitions
RealPDEBench is pinned to revision 0ae4426085d250787f96d5456bc0d30ce5f277d2. Real-
only and simulation-then-real arms share the experimental trajectory budget, test panels,
and scoring rules. Up to 16 training and four validation windows are used per trajectory;
duplicate validation coordinates are counted once, while published test multiplicity is pre-
served. Full test panels and OOD trajectory subsets are distinguished in Section 6.
The scalar development panel uses one Matbench fold and one TDC scaffold seed. The full-
split panel uses five Matbench folds and three TDC scaffold seeds. Reported Matbench and
TDC errors average fold or seed MAEs in native units. Scalar model selection uses validation
MAE and retains the same raw descriptors across baselines and the residual pathway.





\noindent Shodh AI Technical Report\quad 15\par


% Extracted chunk 16


\noindent SHODH AI RESEARCH\quad LUCAN-10B-Physics\par



A.3 Selection, uncertainty and numerical replay
Primary rollout heads are selected on a separate development-validation metric. A v2
residual replaces v1 only when that matched metric improves. The original readout ablation
and the staged rollout scorecard are separate protocols and should not be pooled into one
readout comparison.
The donor-by-calibration factorial reports three-seed arm means and paired 95\% effect in-
tervals. The available summaries do not supply per-arm factorial confidence intervals or
per-row scalar uncertainty, and no such intervals are inferred from rounded point estimates.
Analytical operator comparisons use 64 cases per task, paired TCPI replay uses 16, and the
earlier Stefan timing workload uses a separate 1,000-case cohort.
Reproduction requires the frozen backbone, task encoders and heads, exact partitions, nor-
malization statistics, validation-selection decisions, numerical references, valid-cell masks,
and cell weights. Prediction-level records are needed to recompute paired effects and uncer-
tainty. H100 replay must record dtype, routing precision, token length, batch size, paral-
lelism mode, synchronization, timing boundaries, and allocated memory. This report pro-
vides the numerical tables and protocol distinctions; it does not replace the model weights
and prediction records needed for execution.





\noindent Shodh AI Technical Report\quad 16\par


% Extracted chunk 17


\noindent SHODH AI RESEARCH\quad LUCAN-10B-Physics\par



BIBLIOGRAPHY


\section*{References}


[1] Zongyi Li, Nikola Kovachki, Kamyar Azizzadenesheli, et al. Fourier Neural Operator for
Parametric Partial Differential Equations. ICLR, 2021. https://arxiv.org/abs/2010.08895.
[2] Olaf Ronneberger, Philipp Fischer, and Thomas Brox. U-Net: Convolutional Networks for
Biomedical Image Segmentation. MICCAI, 2015. https://arxiv.org/abs/1505.04597.
[3] Michael McCabe, Bruno Régaldo-Saint Blancard, Liam Holden Parker, et al. Multiple Physics
Pretraining for Physical Surrogate Models. arXiv:2310.02994, 2023. https://arxiv.org/abs/
2310.02994.
[4] Maximilian Herde, Bogdan Raonić, Tobias Rohner, et al. Poseidon: Efficient Foundation
Models for PDEs. NeurIPS, 2024; arXiv:2405.19101. https://arxiv.org/abs/2405.19101.
[5] Michael McCabe, Payel Mukhopadhyay, Tanya Marwah, et al. Walrus: A Cross-Domain
Foundation Model for Continuum Dynamics. arXiv:2511.15684, first submitted 19 November
2025. https://arxiv.org/abs/2511.15684v1.
[6] Tri Dao and Albert Gu. Transformers are SSMs: Generalized Models and Efficient Algorithms
Through Structured State Space Duality. ICML, 2024; arXiv:2405.21060. https://arxiv.
org/abs/2405.21060.
[7] Opher Lieber, Barak Lenz, Hofit Bata, et al. Jamba: A Hybrid Transformer–Mamba Language
Model. arXiv:2403.19887, 2024. https://arxiv.org/abs/2403.19887.
[8] Aran Komatsuzaki, Joan Puigcerver, James Lee-Thorp, et al. Sparse Upcycling: Training
Mixture-of-Experts from Dense Checkpoints. ICLR, 2023; arXiv:2212.05055. https://arxiv.
org/abs/2212.05055.
[9] Taishi Nakamura, Takuya Akiba, Kazuki Fujii, Yusuke Oda, Rio Yokota, and Jun Suzuki.
Drop-Upcycling: Training Sparse Mixture of Experts with Partial Re-initialization. ICLR,
2025; arXiv:2502.19261. https://arxiv.org/abs/2502.19261.
[10] Makoto Takamoto, Timothy Praditia, Raphael Leiteritz, et al. PDEBench: An Exten-
sive Benchmark for Scientific Machine Learning. NeurIPS Datasets and Benchmarks, 2022;
arXiv:2210.07182. https://arxiv.org/abs/2210.07182.
[11] Ruben Ohana, Michael McCabe, Lucas Meyer, et al. The Well: A Large-Scale Collection of
Diverse Physics Simulations for Machine Learning. NeurIPS Datasets and Benchmarks, 2024;
arXiv:2412.00568. https://arxiv.org/abs/2412.00568.
[12] Peiyan Hu, Haodong Feng, Hongyuan Liu, et al. RealPDEBench: A Benchmark for Complex
Physical Systems with Real-World Data. ICLR, 2026. https://openreview.net/forum?id=
y3oHMcoItR.
[13] David Rogers and Mathew Hahn. Extended-Connectivity Fingerprints. Journal of Chemical
Information and Modeling, 50(5):742–754, 2010. https://doi.org/10.1021/ci100050t.
[14] Alexander Dunn, Qi Wang, Alex Ganose, Daniel Dopp, and Anubhav Jain. Benchmarking
Materials Property Prediction Methods: The Matbench Test Set and Automatminer Ref-
erence Algorithm. npj Computational Materials, 6:138, 2020. https://doi.org/10.1038/
s41524-020-00406-3.
[15] Kexin Huang, Tianfan Fu, Wenhao Gao, et al. Therapeutics Data Commons: Machine Learn-
ing Datasets and Tasks for Drug Discovery and Development. NeurIPS Datasets and Bench-
marks, 2021; arXiv:2102.09548. https://arxiv.org/abs/2102.09548.





\noindent Shodh AI Technical Report\quad 17\par


\end{document}
